\documentclass[10pt]{article}
\usepackage[a4paper,margin=19mm]{geometry}
\usepackage[T1]{fontenc}
\usepackage{amsmath,amssymb,amsthm,mathtools,booktabs,array}
\newtheorem{theorem}{Theorem}
\newcommand{\eps}{\varepsilon}
\newcommand{\Du}{D_u}
\newcommand{\Dv}{D_v}
\newcommand{\Q}{\mathbb Q}
\newcommand{\R}{\mathbb R}
\newcommand{\coef}[1]{[#1]}
\title{Fixed testing functions for the strong contact branch}
\author{}
\date{}
\begin{document}
\maketitle
\vspace{-7mm}
The complete state quadratic form can be completed uniformly when the
native contact parameter has a fixed positive lower bound. The construction
below gives its six by six Gram matrix by polynomial division, then assigns
source-only positive matrices to fixed testing functions. It complements
the weak-contact construction in \emph{Completing the state quadratic form
at a singular contact}. All source and testing variables are independent.

\paragraph{Native source.}
Throughout this note \(H,z,y,r,t\) mean the original native source parameters,
rather than the compact variables used for the weak branch. Put
\[
 \begin{gathered}
 \eps=2^{-64},\quad 2^{-12}\le H<1,\quad
 0<z,y\le\eps H,\quad 0<r<1,\quad t>0,\\
 a_c=1+zr,\quad j=a_c-t>0,\quad B=H+yt,\quad S=a_c+yt.
 \end{gathered}
\]
Define
\begin{align}
 \Delta={}&z+(1+z)B+(1+z)B^2+2z(1+z)Br+z^3r^2+2z^2r
       +y\{(1+z)B+z+z^2r(2-r)\},\nonumber\\
 G_s={}&B+z(H+r-1),\qquad
 M_s=(1+z)(B+zr)(H+y)-z(1+y)(1-r)yt,                 \label{eq:sourcepolys}\\
 n={}&\frac{(B+zr)a_cS-r\Delta}{2r\Delta}\ge0,\qquad
 m=\frac{z^2(1+y)(1-r)a_cS}{2G_s\Delta},\nonumber\\
 C={}&\frac{a_cM_s}{2yt\Delta},\qquad
 A=\frac{ya_c(B+zr)S}{2H\Delta},\qquad
 D=\frac{(1+z)(B+zr)(B-zyr)S}{2yj\Delta},\qquad
 \alpha=\frac{(B+zr)M_s}{2(1-H)\Delta}.              \label{eq:source}
\end{align}
The condition \(n\ge0\) is the physical source restriction.
We normalize \(d=1,\lambda=\beta=0\).
For independent real testing angles and state coordinates set
\[
 \Du=1+u^2,\quad \Dv=1+v^2,\quad
 J=uW+uvX-Y-vZ,\qquad M=vW-X+uvY-uZ
\]
and write the complete, unnormalized slack as
\begin{equation}
\begin{split}
 F={}&4n\Du\Dv X^2+4mM^2+4C\Dv(uW-Y)^2+4A\Du\Dv Y^2\\
 &+4DJ^2+4\alpha\Du(vY-Z)^2
       +2\{J^2+\Du(vY-Z)^2+\Dv(X+uZ)^2-\Du\Dv Z^2\}.
\end{split}\label{eq:fullparent}
\end{equation}
Its compact value is \(S_{\mathbf c}=F/(4\Du\Dv)\).

\begin{theorem}\label{thm:strong}
Every source \eqref{eq:source} in the displayed domain has a Gram
representation of \(S_{\mathbf c}\), for all real state coordinates and
\[
 |u|\le2^{-12},\qquad |v|\le2^{-25},
\]
in the following fixed testing family. Its base consists of the twenty
functions \(\mathcal H_{20}/\sqrt{\Du\Dv}\) specified below. For each of the
thirteen nonconstant exponents in Table~\ref{tab:budgets}, append the two
eleven-function blocks
\[
 \frac{\sqrt{1\pm(4096u)^i(33554432v)^j}}{\sqrt{\Du\Dv}}\,
 \mathcal H_{11}.
\]
All Gram entries depend only on the source. The representation includes
literal \(n=0\) and arbitrary nonnegative additions of the eight event
terms.
\end{theorem}

\section{The complete remainder and its polynomial Gram}

Let \(k=\sqrt y\) and define the pure-source quantities
\begin{equation}
\begin{gathered}
 s=a_c/B,\quad c_*=Cyt,\quad d_*=Dyj,\quad
 \ell=y(C+D),\quad q=D/(C+D),\quad g=yCD/(C+D),\\
 \ell_*=H/4,\quad \alpha_*=H^2/8,\quad
 \mu=m/z^2,\quad \nu=nr,\quad a_\alpha=\alpha(1-H),\quad a_A=A/y .
\end{gathered}\label{eq:scaledsource}
\end{equation}
The testing rows are
\begin{align}
 \xi&=X+zrvW/B,& \eta&=Y-HuW/B,& \zeta&=Z-uvW/B,\nonumber\\
 R_h&=(uvX-vZ+syuW)/k,&
 N&=\{Y-quvX+qvZ-(H/B+qsy)uW\}/k,\nonumber\\
 R&=v(u^2-z)W,& T&=uk(u^2-z)W.                     \label{eq:rows}
\end{align}
In particular \(\eta=k(N+qR_h)\).
The four terms
\begin{equation}
 \mathcal R_+=4(\ell-\ell_*)N^2+4Cv^2\eta^2
       +4n\Du\Dv\xi^2+4(\alpha-\alpha_*)\Du(v\eta-\zeta)^2
                                                               \label{eq:retained}
\end{equation}
remain whole. Their coefficients are nonnegative by Section~3.

Here is a polynomial prescription for every entry of the remaining Gram.
In the following expression regard \(N,\xi,\zeta,R_h,W\) as indeterminates,
and abbreviate
\[
 X_0=\xi-zrvW/B,\quad Y_0=k(N+qR_h)+HuW/B,\quad
 Z_0=\zeta+uvW/B,\quad K=ukW,
\]
\[
 J_0=uW-Y_0+v(uX_0-Z_0),\qquad
 M_0=vW-X_0+uvY_0-uZ_0.
\]
The complete remainder is
\begin{equation}
\begin{split}
 \mathcal G={}&4\ell_*N^2+4gR_h^2
 +\frac{8(d_*-c_*)}{B}KN-8gsKR_h
 +\frac{4(c_*t+d_*j)}{B^2}K^2\\
 &+\frac{4c_*t}{B^2}v^2K^2-\frac{8c_*}{B}v^2K(N+qR_h)\\
 &+4\nu\Du\Dv\left(\frac{z^2rv^2}{B^2}W^2-\frac{2zv}{B}W\xi\right)\\
 &+4\alpha_*\Du(vk(N+qR_h)-\zeta)^2\\
 &+4a_\alpha\Du\left\{\frac{(1-H)u^2v^2}{B^2}W^2
       -\frac{2uv}{B}W(vk(N+qR_h)-\zeta)\right\}\\
 &+4z^2\mu M_0^2+4ya_A\Du\Dv Y_0^2
 +2\{J_0^2+\Du(vY_0-Z_0)^2+\Dv(X_0+uZ_0)^2-\Du\Dv Z_0^2\}.
\end{split}\label{eq:G}
\end{equation}
Expanding the event squares gives \(F=\mathcal G+\mathcal R_+\) at
\eqref{eq:rows}. Thus both state normals are still present.

Put \(U=u^2,V=v^2\). Coefficients below refer to the five state
indeterminates in \(\mathcal G\). Define
\begin{align}
 \rho&=\frac{\coef{R_hW}\mathcal G}{uk},&
 p_\xi&=\frac{\coef{\xi W}\mathcal G}{v}+u^2\rho,&
 p_\zeta&=\frac{\coef{\zeta W}\mathcal G}{uv}-\rho,\nonumber\\
 p_N&=\frac{\coef{NW}\mathcal G}{uk},&
 \phi&=\coef{W^2}\mathcal G-su^2(v^2-y)\rho.          \label{eq:coeff}
\end{align}
Each division in \eqref{eq:coeff} is exact polynomial cancellation;
the resulting four expressions \(p_\xi,p_\zeta,p_N,\phi\) are
polynomials in \(U,V\). They define their values also on \(u=0\) or \(v=0\).
Take the following monic divisions, which uniquely define all quotients
and remainders:
\begin{align}
 p_\xi&=(U-z)A_\xi+z\{(V-y)B_\xi+r_\xi\},\nonumber\\
 p_\zeta&=(U-z)A_\zeta+(V-y)B_\zeta+r_\zeta,\qquad
 p_N=(U-z)A_N+(V-y)B_N+r_N,                         \label{eq:divnormal}\\
 d_{\mathrm{rad}}&=\frac4{B^2}\{z^2\mu(B+zr)^2+(\nu r+r^2/2)z^2\},\quad
 e_s=\frac4{B^2}(\nu r+r^2/2)z^2,\nonumber\\
 \phi-\frac{d_{\mathrm{rad}}+e_{\mathrm{rad}}V}{z^2}V(U-z)^2
       &=a_2(V)U^2+a_1(V)U,\nonumber\\
 a_2&=(V-y)A_s+r_a,\qquad a_1=(V-y)B_s+r_b,\qquad
 B_s+zA_s=(V-y)C_s+r_c.                            \label{eq:divscalar}
\end{align}
The symbol \(B_s\) in the last line is a quotient polynomial, unrelated
to the native scalar \(B\). No division by a testing value is used.

To explain why the six remainders vanish, the native stationary contact is
\[
 u=\sqrt z,\quad v=\sqrt y,\quad
 (W,X,Y,Z)=\left(B/\sqrt{zy},-r\sqrt z,H/\sqrt y,1\right).
\]
At this point \(N=\xi=\zeta=R_h=R=T=0\).
The rational weights \eqref{eq:source} give zero parent value and zero
first derivatives in the three independent normal directions and in
\(u,v\). After the \(\rho\)-elimination, the three normal derivatives give
\(p_\xi(z,y)=p_\zeta(z,y)=p_N(z,y)=0\).
The value and two angle derivatives give
\(\phi(z,y)=\phi_U(z,y)=\phi_V(z,y)=0\).
Equations \eqref{eq:divnormal}--\eqref{eq:divscalar} then give
\begin{equation}
 r_\xi=r_\zeta=r_N=r_a=r_b=r_c=0.                    \label{eq:zero}
\end{equation}
Indeed the scalar value and derivatives are respectively
$z^2r_a+zr_b$, $2zr_a+r_b$ and $zr_c$; since $z>0$,
their vanishing determines all three scalar remainders.
These are exact source identities: equivalently, substitute
\eqref{eq:source} in the six explicitly defined remainders and clear
their positive source denominators.

Let \(\mathcal G_0\) be the part of \(\mathcal G\) quadratic in
\(N,\xi,\zeta,R_h\), with no \(W\). In the next formula evaluate every
quotient at \(U=u^2,V=v^2\), and put \(Q_n=u\xi-\zeta\):
\begin{equation}
\begin{split}
 \mathcal O={}&\mathcal G_0+
 \{A_\xi-(v^2-y)B_\xi\}R\xi
 +\frac{B_\xi}{s}(uv^2Q_n-uvkR_h)\xi\\
 &+uA_\zeta R\zeta+\frac{B_\zeta}{s}(v^2Q_n-vkR_h)\zeta
 +A_NTN+\frac{B_N}{s}(vkQ_n-yR_h)N\\
 &+\frac{d_{\mathrm{rad}}+e_{\mathrm{rad}}v^2}{z^2}R^2
 +\frac{A_s}{s}(-R_hT+uRQ_n)
 +\frac{C_s}{s^2}(v^2Q_n^2-2vkR_hQ_n+yR_h^2).
\end{split}\label{eq:O}
\end{equation}
This explicitly defines the symmetric six by six matrix
\[
 f=(N,\xi,\zeta,R_h,R,T)^T,\qquad
 Q=\tfrac12\nabla_f^2\mathcal O,\qquad \mathcal O=f^TQf.
\]
Indeed, using \(kR_h=vQ_n+su(y-v^2)W\), direct monic division gives
\begin{equation}
 \mathcal G-\mathcal O=
 vzr_\xi\xi W+uvr_\zeta\zeta W+ukr_NNW+
 \{u^4r_a+u^2r_b+u^2(v^2-y)r_c\}W^2.               \label{eq:exactdifference}
\end{equation}
Together with \eqref{eq:zero}, this proves the entire parent identity.
In particular \(d_*,t,j\) occur only in the vanishing remainders and
need no interval replacement in \(Q\).

\section{Removing the degenerate contact directions}

The exact testing relation
\begin{equation}
 \mathcal C=T^2-u^2R^2-\frac{u^2-z}{s}R_hT
                 +\frac{u(u^2-z)}sR(u\xi-\zeta)=0             \label{eq:conductor}
\end{equation}
follows at once from \eqref{eq:rows}. Add \(4\mathcal C/H\) to
\(f^TQf\), and use
\begin{equation}
 \widehat f=\left(\sqrt HN,\ \xi+u\zeta+\frac R{H(1+H)},\
 H\zeta,\ \sqrt HR_h,\ R/\sqrt H,\ T/\sqrt H\right)^T. \label{eq:newroots}
\end{equation}
If \(h=\sqrt H\), the matrix taking \(\widehat f\) to \(f\) has diagonal
\((h^{-1},1,H^{-1},h^{-1},h,h)\) and additional entries
\(L_{23}=-u/H,\ L_{25}=-1/[h(1+H)]\).
Thus the exact normalized matrix is
\begin{equation}
 \widehat Q=L^T\left(Q+\frac4H\,
                      \frac{\nabla_f^2\mathcal C}{2}\right)L. \label{eq:normalized}
\end{equation}
Its coefficients are obtained from the displayed polynomial operations;
no entries are chosen by a semidefinite feasibility problem.

The leading source matrix is
\begin{equation}
 Q_*=\operatorname{diag}
 \left(1,2,\tfrac12,\frac2{1+H},\frac2{(1+H)^2},4\right)
 +\frac{e_4e_6^T+e_6e_4^T}{1+H}\succeq\tfrac12 I_6. \label{eq:baseline}
\end{equation}
For the last inequality, the fifth diagonal after subtraction of
\(I_6/2\) is \((1-H)(3+H)/[2(1+H)^2]\). The fourth/sixth block has
positive first pivot and determinant
\([17+7H(2-H)]/[4(1+H)^2]\); the other entries are nonnegative.

\section{Uniform source and coefficient bounds}

Set \(a_0=1/[2(1+H)]\). The nine errors are defined by
\begin{equation}
\begin{array}{lll}
 B/H=1+e_B,&sH=1+e_s,&c_*/H=a_0+e_C,\\
 g/H=a_0+e_g,&a_\alpha/H^2=a_0+e_\alpha,&AH/y=a_0+e_A,\\
 mH^2/z^2=(1-r)a_0+e_m,&nr=a_0-r/2+e_n,&q=t+e_q.
\end{array}\label{eq:errors}
\end{equation}
They satisfy
\begin{equation}
 |e_B|\le2\eps,\quad |e_s|\le4\eps,\quad
 |e_C|,|e_\alpha|,|e_A|,|e_m|,|e_n|\le16\eps,\quad
 |e_g|,|e_q|\le64\eps.                              \label{eq:errorbounds}
\end{equation}
Here is a factorwise proof uniform in \(r,t,j,1-H\).
Writing \(T_s=B+zr\) and \(U_s=B-zyr\), the source formulas give
\begin{equation}
\begin{gathered}
 1\le T_s/H,\ S\le(1+\eps)^2,\qquad
 1-\eps^2\le U_s/H\le(1+\eps)^2,\\
 1-\eps\le G_s/H\le1+3\eps,\quad
 1\le M_s/H^2\le(1+\eps)^4,\quad
 1\le\Delta/[H(1+H)]\le1+10\eps .                  \label{eq:envelopes}
\end{gathered}
\end{equation}
The lower bound for \(M_s\) uses
\[
 M_s=(1+z)T_sH+yG_s+yzj+zr(yt)(1+y);
\]
its upper bound uses its first expression in \eqref{eq:sourcepolys}
and drops the nonnegative subtracted term. The other elementary bounds
follow from \(t<a_c\le1+\eps\).
For completeness, \(B\le q_0H\) with \(q_0=1+\eps+\eps^2\).
Termwise bounds in \(\Delta\) give
\[
 \Delta-H(1+H)\le H(a+bH+cH^2),\quad
 a=2\eps+\eps^2\le3\eps,
\]
\[
 b=6\eps+10\eps^2+6\eps^3+\eps^4\le7\eps,\quad
 c=\eps+5\eps^2+8\eps^3+5\eps^4+\eps^5\le2\eps .
\]
Since \(3+7H+2H^2\le6(1+H)\) for \(0\le H\le1\), this even
gives the stronger upper bound \(1+6\eps\).

The ratios of \(c_*/H,d_*/H,a_\alpha/H^2,AH/y\) to \(a_0\)
have lower bound \((1-\eps^2)/(1+10\eps)\) and respective upper
bounds \((1+\eps)^5,(1+\eps)^7,(1+\eps)^6,(1+\eps)^5\).
Their relative errors are at most \(12\eps\).
The ratio of \(mH^2/z^2\) to \((1-r)a_0\) lies between
\[
 \frac1{(1+3\eps)(1+10\eps)}
 \quad\hbox{and}\quad \frac{(1+\eps)^4}{1-\eps},
\]
and \(nr+r/2=AH/y\) exactly. These imply the corresponding
\(16\eps\) bounds. Also
\(|e_B|\le\eps(1+\eps)\) and \(|e_s|\le3\eps\).
If \(c_*/H=a_0(1+\delta_C)\), \(d_*/H=a_0(1+\delta_D)\),
put \(\delta_w=(j\delta_C+t\delta_D)/a_c\).
Then \(|\delta_w|\le12\eps\) and
\[
 q=\frac t{a_c}\frac{1+\delta_D}{1+\delta_w},\qquad
 \frac gH=\frac{a_0}{a_c}
           \frac{(1+\delta_C)(1+\delta_D)}{1+\delta_w}.
\]
The endpoint ratios give \(|q-t|<26\eps\) and
\(|g/H-a_0|<37\eps/2\), proving \eqref{eq:errorbounds}.
In particular
\[
 \ell\ge4g\ge H(1-256\eps)>H/4,\qquad
 \alpha\ge H^2(1/4-16\eps)>H^2/8.
\]
This proves positivity of every coefficient in \eqref{eq:retained}.

The following finite majorant specifies the coefficient estimates
without long rational-entry lists. In \eqref{eq:normalized} substitute
\[
 H=h^2,\quad B=h^2(1+e_B),\quad s=(1+e_s)/h^2,\quad
 z=h^2z_1,\quad k=hk_1,
\]
and replace the other quantities by \eqref{eq:errors}, with
\(\ell_*=h^2/4,\alpha_*=h^4/8\). Expand
\(\widehat Q=\sum C_{ij}u^iv^j\).
Absorb any positive rational constant in a denominator into its numerator.
Every remaining denominator factors into
\(h^a(1+h^2)^b(1+e_B)^c(1+e_s)^d\).
For a numerator monomial with \(h\)-exponent \(a'\), bound its \(h\)
factor by \(64^{\max(a-a',0)}\), its remaining source factors using
\eqref{eq:errorbounds},
\[
 |r|\le1,\quad |t|\le2,\quad |z_1|\le\eps,\quad
 |k_1|\le2^{-32},
\]
and its denominator by \(2^{c+d}\). Sum absolute coefficients and
count symmetric off-diagonal entries twice. The triangle inequality
bounds the operator norm by this explicitly defined finite sum.
The result for \(C_{00}-Q_*\), and for each
\(2^{-12i-25j}C_{ij}\), is bounded by Table~\ref{tab:budgets}.
It uses only the divisions \eqref{eq:divnormal}--\eqref{eq:divscalar}
and the two Hessians above; thus it also specifies every rational
coefficient estimate.
\begin{table}[ht]
\centering\small
\begin{tabular}{c*{7}{c}}
\toprule
\((i,j)\)&00&01&02&10&11&12&20\\
\(\alpha_{ij}\)&\(2^{-40}\)&\(2^{-50}\)&\(2^{-33}\)&\(2^{-4}\)&\(2^{-55}\)&\(2^{-41}\)&\(2^{-18}\)\\
\midrule
\((i,j)\)&21&22&30&32&40&42&50\\
\(\alpha_{ij}\)&\(2^{-62}\)&\(2^{-48}\)&\(2^{-28}\)&\(2^{-65}\)&\(2^{-42}\)&\(2^{-72}\)&\(2^{-51}\)\\
\bottomrule
\end{tabular}
\caption{Fixed rational coefficient budgets. The constant entry bounds
\(\|C_{00}-Q_*\|_2\); all others bound
\(\|2^{-12i-25j}C_{ij}\|_2\).}\label{tab:budgets}
\end{table}
There are exactly fourteen monomials and fifty-one nonzero
upper-triangular coefficient records. Their full exact definition is
\[
 (i,j,a,b)\longmapsto
 [u^av^b]\left\{
 L^T\left(\tfrac12\nabla_f^2\mathcal O
        +\tfrac4H\,\tfrac12\nabla_f^2\mathcal C\right)L
 \right\}_{ij},
\]
after the displayed source substitution. Off-diagonal entries are
actual symmetric entries.

\section{Positive fixed Gram blocks and state restoration}

The sum of all fourteen budgets is
\[
 \frac{295165937727343559809}{4722366482869645213696}<\frac18.
\]
For each nonconstant exponent set
\[
 p_{ij}=(4096u)^i(33554432v)^j,\qquad
 A_{ij}=2^{-12i-25j}C_{ij},\qquad
 Q_{\rm res}=C_{00}-\sum_{(i,j)\ne(0,0)}\alpha_{ij}I_6 .
\]
Then \(Q_{\rm res}\succeq3I_6/8\),
\((\alpha_{ij}I_6\pm A_{ij})/2\succeq0\), and the exact identity
\[
 \widehat Q=Q_{\rm res}+
 \sum_{(i,j)\ne(0,0)}
 \left\{\frac{\alpha_{ij}I_6+A_{ij}}2(1+p_{ij})
       +\frac{\alpha_{ij}I_6-A_{ij}}2(1-p_{ij})\right\}
\]
proves positivity on the closed testing rectangle.

All six roots \eqref{eq:newroots} are linear combinations, with source-only
coefficients, of
\[
 \mathcal H_{11}=(X,Y,Z,vW,uvW,uW,uvX,vZ,u^2vW,u^3W,uZ).
\]
The retained terms \eqref{eq:retained} use this list together with
\[
 (uX,vX,v^2W,uv^2W,vY,uvY,W,uvZ,uY),
\]
giving \(\mathcal H_{20}\). For example the \(n\)-term uses
\(\xi,u\xi,v\xi,uv\xi\), and the last retained term uses
\(v\eta-\zeta,u(v\eta-\zeta)\). The source expansions are polynomial
in the original state coordinates, so arbitrary state signs and
zero state coordinates are included.

Push the six-dimensional positive matrices to these fixed spans,
restore \(\mathcal R_+\), divide the Gram coefficients by four and
the testing functions by \(\sqrt{\Du\Dv}\).
This gives precisely \(S_{\mathbf c}\), proving Theorem~\ref{thm:strong}.
No formula divides by \(n\), so \(n=0\) is literal.
The event roots already belong to \(\mathcal H_{20}\); their nonnegative
addition preserves the representation.

Together with the weak-contact family and the two preserved remote
testing matrices, this supplies the native source envelope with one
fixed certificate of total PSD order \(635\), including nine source
signs. The remote matrices provide the projective testing endpoints.
Admission of every canonical source is a separate requirement; the
present theorem establishes no equality with the full canonical cone.
\end{document}
